% Section 3 — to be inserted after Section 2 ("Kobayashi and horizontal geodesics").
% Requires \usepackage{booktabs} and \usepackage{graphicx} (for \resizebox) in the preamble.
% Table data: output/dF_convergence_table.{csv,json,tex}, regenerated by
%   python -m scripts.run_dF_convergence_table

\section{Computing $d_F$ with the BFGS algorithm}

Throughout this section the pair is normalized as in Section~2, $P_0=(0,I_n)$ and $P_1=(\mu,\Sigma)$,
and we compute
\[
d_F(P_0,P_1)^2=\min_{S\in\mathrm{AntiSym}(n)}\tfrac14\,\|\log G(P_1,S)\|_F^2 .
\]

\subsection{A reduced objective}

Evaluating $\|\log G(P_1,S)\|_F$ directly requires the logarithm of a $(2n+1)\times(2n+1)$ matrix.
In practice we use an equivalent objective that only involves an $n\times n$ singular value decomposition.
For $K\in\mathrm{AntiSym}(n)$, set
\[
Y(K)=\frac12\Big(I_n+\Sigma+\frac12\mu\mu^\top\Big)+K,
\qquad
c_i(K)=\sigma_i\big(\Sigma^{-1/2}Y(K)\big),\quad i=1,\dots,n,
\]
where $\sigma_i$ denotes the $i$-th singular value, and
\begin{equation}
J(K)=2\sum_{i=1}^n \operatorname{arcosh}^2\big(c_i(K)\big).
\label{eq:J}
\end{equation}
One checks that $c_i(K)\ge 1$, so that \eqref{eq:J} is well defined, using the identity
$Y^\top\Sigma^{-1}Y-I=(Y-\Sigma)^\top\Sigma^{-1}(Y-\Sigma)+\frac12\mu\mu^\top\succeq0$.

\paragraph{Link with the lift.}
Since $JGJ=G^{-1}$ and $J$ is an orthogonal involution, the spectrum of $G=G(P_1,S)$ is invariant under
$\lambda\mapsto\lambda^{-1}$; its eigenvalues are $1,\lambda_1^{\pm1},\dots,\lambda_n^{\pm1}$ and
$\tfrac14\|\log G\|_F^2=\tfrac12\sum_i\log^2\lambda_i$.
The reduced objective corresponds to $c_i=\cosh\big(\tfrac12\log\lambda_i\big)$, i.e.
\begin{equation}
\tfrac14\,\|\log G(P_1,S)\|_F^2 = J\big(\tfrac12 S\big).
\label{eq:J-lift}
\end{equation}
\emph{Status.} We have verified \eqref{eq:J-lift} numerically on random pairs in dimensions $2$ to $10$
(relative agreement $10^{-10}$ or better); a complete proof is still to be written.
On the example of Section~2 with $S=0$, $\lambda=65.420223$ gives
$\cosh(\tfrac12\log\lambda)=4.1058=c_1(0)$.
Assuming \eqref{eq:J-lift},
\[
d_F^2=\min_{K\in\mathrm{AntiSym}(n)}J(K),\qquad d_0^2=J(0).
\]

We parametrize $K$ by its upper-triangular entries
$x=(K_{ij})_{i<j}\in\mathbb R^{p}$, $p=n(n-1)/2$, and minimize $x\mapsto J(K(x))$ over $\mathbb R^p$.

\paragraph{Normalization in practice.}
Instead of $\Sigma_0^{-1/2}$ we use the Cholesky factorization $\Sigma_0=LL^\top$ and set
$\mu=L^{-1}(\mu_1-\mu_0)$, $\Sigma=L^{-1}\Sigma_1L^{-\top}$, which is allowed by affine invariance.
Similarly, if $\Sigma=L_\Sigma L_\Sigma^\top$ and $W=L_\Sigma^{-1}$, then $W=Q\Sigma^{-1/2}$ with $Q$ orthogonal,
so $W\,Y(K)$ and $\Sigma^{-1/2}Y(K)$ have the same singular values.
The matrices $W$ and $Y(0)$ are computed once per pair; each evaluation of $J$ then costs
one $n\times n$ SVD.

\subsection{Gradient}

Write $\varphi(c)=2\operatorname{arcosh}^2(c)$, so that $J=\sum_i\varphi(c_i)$ and
\[
\varphi'(c)=\frac{4\operatorname{arcosh}(c)}{\sqrt{c^2-1}}\quad(c>1),\qquad \varphi'(1)=\lim_{c\to1}\varphi'(c)=4 .
\]
Let $M(K)=W\,Y(K)=U\operatorname{diag}(c)\,V^\top$ be an SVD with columns $u_i$, $v_i$.
For a simple singular value, $\mathrm dc_i=u_i^\top(\mathrm dM)\,v_i$, and since $\mathrm dM=W\,\mathrm dK$,
\[
\mathrm dJ=\sum_i\varphi'(c_i)\,u_i^\top W\,\mathrm dK\,v_i=\big\langle \Gamma(K),\mathrm dK\big\rangle_F,
\qquad
\Gamma(K)=W^\top U\operatorname{diag}\big(\varphi'(c)\big)V^\top .
\]
In the coordinates $x$, $\mathrm dK=\sum_{i<j}\mathrm dx_{ij}(E_{ij}-E_{ji})$, hence
\begin{equation}
\frac{\partial J}{\partial x_{ij}}=\Gamma_{ij}-\Gamma_{ji},\qquad i<j .
\label{eq:gradJ}
\end{equation}
The same SVD gives $J$ and $\nabla J$, so a gradient costs no more than a function value.
Formula~\eqref{eq:gradJ} is checked against centered finite differences in the test suite.

\begin{proposition}[A critical point at $K=0$]
If $\mu=0$ or $\mu$ is an eigenvector of $\Sigma$, then $\nabla J(0)=0$.
\end{proposition}

\begin{proof}
Under either assumption $\Sigma$ commutes with $\mu\mu^\top$, hence with $Y(0)$.
Then $N=\Sigma^{-1/2}Y(0)$ is symmetric positive definite, its SVD is an eigendecomposition
$N=U\operatorname{diag}(c)U^\top$, and (using $\Sigma^{-1/2}$ in place of $W$, which does not change $\mathrm dJ$)
$\Gamma(0)=\Sigma^{-1/2}\varphi'(N)$.
Both factors are functions of commuting symmetric matrices, so $\Gamma(0)$ is symmetric and
\eqref{eq:gradJ} vanishes.
\end{proof}

This covers equal means ($\mu=0$), equal covariances ($\Sigma=I_n$ after normalization),
isotropic covariances and means aligned with an eigenvector of $\Sigma$.
The proposition only says that $K=0$ is a critical point; that it is the global minimum is proved for
equal means and equal covariances (Proposition~7) and observed numerically in the other cases below.

\subsection{The BFGS iteration}

We minimize $J$ with the quasi-Newton BFGS method
(\texttt{scipy.optimize.minimize}, \texttt{method="BFGS"}) using the analytic gradient \eqref{eq:gradJ}.
The algorithm maintains an approximation $H_k\in\mathrm{SPD}(p)$ of the inverse Hessian of $J$.
It is started at
\[
x_0=0\quad(\text{i.e. }K_0=0,\ J(K_0)=d_0^2),\qquad H_0=I_p,
\]
and at each iteration $k$ it performs:
\begin{enumerate}
\item \textbf{Gradient.} Compute $g_k=\nabla J(x_k)$ from one SVD of $W\,Y(K_k)$.
\item \textbf{Stopping test.} Stop if $\|g_k\|_\infty\le\varepsilon$, with $\varepsilon=10^{-10}$.
\item \textbf{Direction.} $p_k=-H_kg_k$, a descent direction since $H_k$ is positive definite.
\item \textbf{Line search.} Choose a step $\alpha_k>0$ satisfying the Wolfe conditions
\[
J(x_k+\alpha_kp_k)\le J(x_k)+c_1\alpha_k\,g_k^\top p_k,\qquad
\nabla J(x_k+\alpha_kp_k)^\top p_k\ge c_2\,g_k^\top p_k,
\]
with $c_1=10^{-4}$ and $c_2=0.9$.
\item \textbf{Update.} $x_{k+1}=x_k+\alpha_kp_k$, $s_k=x_{k+1}-x_k$, $y_k=g_{k+1}-g_k$, $\rho_k=1/(y_k^\top s_k)$ and
\[
H_{k+1}=\big(I-\rho_ks_ky_k^\top\big)H_k\big(I-\rho_ky_ks_k^\top\big)+\rho_ks_ks_k^\top .
\]
The curvature condition of the line search ensures $y_k^\top s_k>0$, so $H_{k+1}$ stays positive definite.
\end{enumerate}
Each iteration therefore costs one or a few SVDs of size $n\times n$ ($O(n^3)$) plus the rank-two update
of $H_k$ ($O(p^2)=O(n^4)$).
Because every accepted step decreases $J$ and $x_0=0$, the computed value always satisfies $d_F\le d_0$.

Two remarks on reliability.
First, BFGS returns a critical point, not a certified global minimum.
To test this, every case of Table~\ref{tab:bfgs} was also solved from $4$ random starting points
$K_{ij}\sim\mathcal N\big(0,\|Y(0)\|_F^2/n\big)$; all runs agreed with the run started at $K=0$ to within
$10^{-14}$, and none found a lower value.
Second, when $\varepsilon=10^{-10}$ is below the accuracy attainable in floating point, the line search
stops making progress and SciPy reports a \emph{precision loss}; these cases are marked with $\dagger$.
The final gradients are then at most $1.5\times10^{-7}$ and the restarts above still agree to $10^{-14}$.

\subsection{Numerical results}

Table~\ref{tab:bfgs} reports $d_F$ and $d_0$ on fifteen pairs: exact cases, structured cases,
random pairs (both covariances drawn with a prescribed condition number and a random rotation, and the mean
difference drawn in a random direction with prescribed normalized norm), and the stress family
$\Sigma=\operatorname{diag}(1/L,L)$, $\mu=(1,\sqrt L)$.
The columns $\kappa(\Sigma)$ and $\|\mu\|$ refer to the normalized pair, and the relative error is
$|d_0-d_F|/d_F$. ``iter.'' is the number of BFGS iterations and ``eval.'' the number of evaluations of
$(J,\nabla J)$; the time is the best of 5 runs of the full computation of $d_F$ in Python/NumPy.

\begin{table}[h]
\centering
\resizebox{\textwidth}{!}{%
\begin{tabular}{lrrrrrrrrrr}
\toprule
Case & $n$ & $\kappa(\Sigma)$ & $\|\mu\|$ & $d_0$ & $d_F$ & rel.\ err. & iter. & eval. & $\|\nabla J\|_\infty$ & time (ms) \\
\midrule
equal means            & 3  & 16                  & 0    & 1.386294  & 1.386294  & 0      & 0   & 1   & $0$                    & 0.7 \\
equal covariances      & 3  & 1                   & 3.30 & 2.810124  & 2.810124  & 0      & 0   & 1   & $6\cdot10^{-16}$       & 0.6 \\
example of Section 2   & 2  & 20                  & 4.47 & 3.178995  & 3.149173  & 0.947\%  & 5   & 6   & $4\cdot10^{-14}$       & 2.7 \\
isotropic              & 3  & 1                   & 3.46 & 2.671396  & 2.671396  & 0      & 0   & 1   & $5\cdot10^{-16}$       & 0.5 \\
diagonal, aligned mean & 3  & 100                 & 3.00 & 2.724948  & 2.724948  & 0      & 0   & 1   & $0$                    & 0.5 \\
diagonal, mixed mean   & 3  & 100                 & 3.46 & 4.059695  & 3.997851  & 1.547\%  & 8   & 11  & $2\cdot10^{-12}$       & 5.7 \\
dense, unnormalized    & 3  & 52                  & 3.22 & 2.967103  & 2.925964  & 1.406\%  & 8   & 11  & $6\cdot10^{-12}$       & 5.5 \\
random                 & 5  & 9.0                 & 1    & 1.719896  & 1.719578  & 0.018\%  & 14  & 25  & $4\cdot10^{-15}$       & 10.9 \\
random                 & 5  & $2.9\cdot10^{3}$    & 5    & 5.718017  & 5.687416  & 0.538\%  & 20  & 59  & $8\cdot10^{-9}\,\dagger$  & 27.9 \\
random, large mean     & 5  & 45                  & 30   & 10.051179 & 10.049849 & 0.013\%  & 32  & 69  & $2\cdot10^{-9}\,\dagger$  & 38.5 \\
random                 & 10 & $2.3\cdot10^{3}$    & 3    & 5.372293  & 5.310237  & 1.169\%  & 52  & 106 & $9\cdot10^{-10}\,\dagger$ & 57.6 \\
random                 & 20 & $3.5\cdot10^{5}$    & 10   & 9.800225  & 9.646316  & 1.596\%  & 207 & 227 & $6\cdot10^{-8}\,\dagger$  & 720 \\
random                 & 30 & 48                  & 3    & 4.411788  & 4.397404  & 0.327\%  & 52  & 69  & $2\cdot10^{-7}\,\dagger$  & 748 \\
stress $L=100$         & 2  & $10^{4}$            & 10   & 5.912835  & 5.413124  & 9.231\%  & 7   & 44  & $3\cdot10^{-9}\,\dagger$  & 24.4 \\
stress $L=1000$        & 2  & $10^{6}$            & 31.6 & 8.993768  & 7.717864  & 16.532\% & 5   & 10  & $2\cdot10^{-12}$       & 4.6 \\
\bottomrule
\end{tabular}}
\caption{BFGS computation of $d_F$ started at $K=0$, compared with $d_0$.
$\dagger$: SciPy stopped on precision loss before reaching $\|\nabla J\|_\infty\le10^{-10}$
(see the remarks above).}
\label{tab:bfgs}
\end{table}

\paragraph{Observations.} These are numerical observations on the fifteen pairs of
Table~\ref{tab:bfgs}, not general statements.
\begin{itemize}
\item In the four cases covered by the proposition above, BFGS stops at iteration $0$ and $d_F=d_0$.
For equal means and equal covariances the result matches the closed forms (8) and (9)
to $2\cdot10^{-16}$.
\item In dimensions $2$ and $3$, BFGS converges in $5$ to $8$ iterations. The number of iterations grows with
the dimension and the conditioning: $14$--$32$ for $n=5$, $52$ for $n=10$ and $n=30$, and $207$ for $n=20$
with $\kappa(\Sigma)\approx3.5\cdot10^5$.
\item Computing $d_F$ takes from a few milliseconds ($n\le3$) to less than one second ($n=30$),
whereas $d_0$ takes about $0.5$\,ms in all cases.
\item Outside the stress family, the relative error of $d_0$ stays below $1.6\%$. It reaches $9.2\%$ and $16.5\%$
in the stress cases, where the covariance is very anisotropic and the mean mixes its two eigendirections.
\end{itemize}
